% !PRINTED none
% Abschnitt-opening page: no running head/page number printed on this recto.
\begin{center}
\gesperrt{Twelfth Chapter.}

\textbf{The Double Tangents of a Curve of the Fourth Order.}

\rule{0.3\linewidth}{0.4pt}

§. 95.

\gesperrt{Number of the Double Tangents of a Curve\\
of the Fourth Order.}
\end{center}

A geometrical problem which, on account of its manifold relations to other regions, is of particular importance, and which at the same time, in a similar way to the problem of the points of inflection of the curves of the third order, offers remarkable algebraic relations, is that of the double tangents of the curves of the fourth order.

By a \gesperrt{double tangent} of a curve one understands, as has already been remarked in \S.\,86, a straight line which touches the curve in two distinct points. In curves of lower than the fourth order, double tangents cannot occur.

In the curves of the fourth order a double tangent has, besides the points of contact, no point in common with the curve. In special cases the two points of contact may also coincide. Then we have lines with four-point contact.

The determination of the double tangents will, as an algebraic problem, depend on a certain algebraic equation whose degree is equal to the number of the double tangents, and the first question is that of the degree of this equation, that is, of the number of the double tangents. We restrict ourselves here to the consideration of curves of the fourth order, although the way that we take is applicable also to curves of higher order
%CONT>>
